\documentclass[UTF8,fontset=fandol,aspectratio=169,11pt]{ctexbeamer}
% Shared Beamer preamble for Big Data and Management Decision-Making slides.

\usetheme{Madrid}
\usecolortheme{default}
\usefonttheme{professionalfonts}

\usepackage{amsmath,amssymb,mathtools}
\usepackage{booktabs,array}
\usepackage{tikz}
\usepackage{graphicx}
\usepackage{listings}
\usepackage{xcolor}
\usetikzlibrary{arrows.meta,positioning,shapes.geometric}

\newcommand{\BDMFandolPath}{/usr/local/texlive/2025/texmf-dist/fonts/opentype/public/fandol/}
\setCJKmainfont[
  Path=\BDMFandolPath,
  UprightFont=FandolSong-Regular.otf,
  BoldFont=FandolSong-Regular.otf,
  ItalicFont=FandolKai-Regular.otf,
  BoldItalicFont=FandolKai-Regular.otf
]{FandolSong-Regular.otf}
\setCJKsansfont[
  Path=\BDMFandolPath,
  UprightFont=FandolSong-Regular.otf,
  BoldFont=FandolSong-Regular.otf
]{FandolSong-Regular.otf}
\setCJKmonofont[
  Path=\BDMFandolPath,
  UprightFont=FandolSong-Regular.otf,
  BoldFont=FandolSong-Regular.otf
]{FandolSong-Regular.otf}
\setmonofont{Menlo}

\definecolor{AcademicBlue}{RGB}{22,44,73}
\definecolor{AcademicTeal}{RGB}{22,119,119}
\definecolor{AcademicGray}{RGB}{76,84,95}
\definecolor{AcademicLight}{RGB}{245,247,250}
\definecolor{AcademicAmber}{RGB}{181,111,35}
\definecolor{CodeBg}{RGB}{248,248,248}

\setbeamercolor{structure}{fg=AcademicBlue}
\setbeamercolor{frametitle}{fg=white,bg=AcademicBlue}
\setbeamercolor{title}{fg=white,bg=AcademicBlue}
\setbeamercolor{block title}{fg=white,bg=AcademicBlue}
\setbeamercolor{block body}{bg=AcademicLight,fg=black}
\setbeamertemplate{navigation symbols}{}
\setbeamertemplate{footline}[frame number]
\setbeamertemplate{itemize item}{\raisebox{0.2ex}{\(\blacktriangleright\)}}
\setbeamerfont{title}{series=\mdseries}
\setbeamerfont{subtitle}{series=\mdseries}
\setbeamerfont{frametitle}{series=\mdseries}
\setbeamerfont{block title}{series=\mdseries}
\setbeamerfont{section in toc}{series=\mdseries}

\lstdefinestyle{pythonstyle}{
  basicstyle=\ttfamily\scriptsize,
  backgroundcolor=\color{CodeBg},
  keywordstyle=\color{AcademicBlue},
  commentstyle=\color{AcademicGray},
  stringstyle=\color{AcademicAmber},
  showstringspaces=false,
  breaklines=true,
  frame=single,
  rulecolor=\color{black!20},
  columns=fullflexible
}

\lstdefinestyle{sqlstyle}{
  language=SQL,
  basicstyle=\ttfamily\scriptsize,
  backgroundcolor=\color{CodeBg},
  keywordstyle=\color{AcademicBlue},
  commentstyle=\color{AcademicGray},
  stringstyle=\color{AcademicAmber},
  showstringspaces=false,
  breaklines=true,
  frame=single,
  rulecolor=\color{black!20},
  columns=fullflexible
}

\lstset{style=pythonstyle}

\hypersetup{
  colorlinks=true,
  linkcolor=AcademicBlue,
  urlcolor=AcademicTeal,
  pdfauthor={大数据与管理决策基础},
  pdfsubject={大数据与管理决策基础课程课件}
}


\title[因果推断与管理干预]{第8讲：因果推断与管理干预}
\subtitle{从预测风险到评估行动是否改变结果}
\author{《大数据与管理决策基础》}
\date{}

\begin{document}

\begin{frame}
  \titlepage
  \begin{center}
    \small 核心命题：预测告诉我们谁更可能出问题，因果推断告诉我们行动是否真的改变结果。
  \end{center}
\end{frame}

\begin{frame}{本讲学习目标}
\begin{enumerate}
  \item 区分相关、预测和因果；
  \item 用潜在结果表达管理干预效果；
  \item 理解反事实、混杂和选择偏差；
  \item 说明随机实验为何是识别基准；
  \item 掌握倾向得分、分层和逆概率加权的直觉；
  \item 使用 DID 评估前后政策变化；
  \item 将因果估计写成带假设的管理建议。
\end{enumerate}
\end{frame}

\begin{frame}{课程主线中的第8周}
\begin{center}
\resizebox{0.96\linewidth}{!}{%
\begin{tikzpicture}[
  node/.style={draw=AcademicBlue!65, fill=AcademicLight, rounded corners=1pt,
               minimum height=0.85cm, minimum width=2.1cm, align=center},
  active/.style={draw=AcademicBlue, fill=AcademicBlue, text=white,
               minimum height=0.9cm, minimum width=2.25cm, align=center},
  arrow/.style={-{Stealth[length=2mm]}, thick, AcademicTeal}
]
\node[node] (kpi) {可信指标};
\node[node, right=0.55cm of kpi] (ab) {随机实验};
\node[node, right=0.55cm of ab] (pred) {预测模型};
\node[active, right=0.55cm of pred] (causal) {因果推断};
\node[node, right=0.55cm of causal] (opt) {优化决策};
\draw[arrow] (kpi) -- (ab);
\draw[arrow] (ab) -- (pred);
\draw[arrow] (pred) -- (causal);
\draw[arrow] (causal) -- (opt);
\end{tikzpicture}
}
\end{center}
\vspace{0.2cm}
\small
第7周回答“谁更可能流失”，第8周回答“对谁采取行动是否真的有效”。
\end{frame}

\begin{frame}{预测问题与因果问题}
\begin{block}{预测}
\[
  \widehat{\mathbb P}(Y=1\mid X=x)
\]
在当前状态下，客户有多大可能流失？
\end{block}
\begin{block}{因果}
\[
  Y(1)-Y(0)
\]
如果发放优惠券，相比不发券，结果是否改变？
\end{block}
\small
高风险客户更可能收到优惠券，也更可能流失；因此处理组结果不能直接解释为处理效果。
\end{frame}

\begin{frame}{潜在结果}
对个体 \(i\)：
\[
  Y_i(1)=\text{接受处理时的结果},
  \qquad
  Y_i(0)=\text{不接受处理时的结果}.
\]
\[
  Y_i=D_iY_i(1)+(1-D_i)Y_i(0).
\]
\[
  ATE=\mathbb E[Y(1)-Y(0)],
  \qquad
  ATT=\mathbb E[Y(1)-Y(0)\mid D=1].
\]
\begin{block}{基本因果问题}
同一个体不能同时观察 \(Y_i(1)\) 和 \(Y_i(0)\)，其中一个永远是反事实。
\end{block}
\end{frame}

\begin{frame}{相关、预测、因果}
\begin{center}
\begin{tabular}{p{0.18\linewidth}p{0.36\linewidth}p{0.32\linewidth}}
\toprule
任务 & 典型问题 & 主要风险\\
\midrule
相关 & 工单数是否与流失同向变化？ & 把共同变化当作原因。\\
预测 & 谁未来最可能流失？ & 把特征解释成干预对象。\\
因果 & 优惠券是否提高留存？ & 忽视混杂和识别假设。\\
\bottomrule
\end{tabular}
\end{center}
\end{frame}

\begin{frame}{随机实验是识别基准}
若处理分配随机：
\[
  D_i\perp (Y_i(1),Y_i(0)).
\]
则
\[
  \mathbb E[Y\mid D=1]-\mathbb E[Y\mid D=0]
  =
  \mathbb E[Y(1)-Y(0)].
\]
\begin{block}{管理含义}
随机化让处理组和对照组在设计上可比；观测性方法必须用假设和调整努力恢复这种可比性。
\end{block}
\end{frame}

\begin{frame}{观测性因果识别假设}
\begin{itemize}
  \item 一致性：观察结果等于实际接受处理版本下的潜在结果；
  \item SUTVA：个体之间没有处理溢出，处理没有隐藏版本；
  \item 条件可交换性：\((Y(1),Y(0))\perp D\mid X\)；
  \item 重叠性：\(0<\mathbb P(D=1\mid X=x)<1\)。
\end{itemize}
\begin{block}{课堂提醒}
这些假设不能由软件自动保证，必须结合业务机制判断。
\end{block}
\end{frame}

\begin{frame}{混杂与选择偏差}
\begin{center}
\resizebox{0.72\linewidth}{!}{%
\begin{tikzpicture}[
  node/.style={draw=AcademicBlue!65, fill=AcademicLight, rounded corners=1pt,
               minimum height=0.9cm, minimum width=2.6cm, align=center},
  arrow/.style={-{Stealth[length=2mm]}, thick, AcademicTeal}
]
\node[node] (risk) {客户风险};
\node[node, right=2.2cm of risk] (coupon) {是否发券};
\node[node, below=1.5cm of coupon] (retention) {是否留存};
\draw[arrow] (risk) -- (coupon);
\draw[arrow] (risk) -- (retention);
\draw[arrow] (coupon) -- (retention);
\end{tikzpicture}
}
\end{center}
\small
若高风险客户更容易被发券，朴素比较同时包含优惠券效果和原始风险差异。
\end{frame}

\begin{frame}{倾向得分}
\[
  e(X)=\mathbb P(D=1\mid X).
\]
\begin{itemize}
  \item 分层：在相近风险或相近倾向得分内比较；
  \item 匹配：寻找处理组和对照组中的相近对象；
  \item 加权：处理组权重 \(1/e(X)\)，对照组权重 \(1/(1-e(X))\)；
  \item 平衡检查：调整后协变量差异是否降低。
\end{itemize}
\[
  \widehat{ATE}_{IPW}
  =
  \frac{1}{n}\sum_i
  \left[
    \frac{D_iY_i}{\hat e(X_i)}
    -
    \frac{(1-D_i)Y_i}{1-\hat e(X_i)}
  \right].
\]
\end{frame}

\begin{frame}{第8周观测性优惠券数据}
\begin{center}
\begin{tabular}{lrrr}
\toprule
方法 & 处理组均值 & 对照组均值 & 效果\\
\midrule
朴素差异 & 0.7500 & 0.7500 & 0.0000\\
风险分层 & 0.7500 & 0.7500 & 0.2500\\
IPW & NA & NA & 0.2199\\
\bottomrule
\end{tabular}
\end{center}
\begin{block}{解释}
朴素差异为 0，不代表优惠券无效；处理组中高风险客户更多，原始比较被选择偏差抵消。
\end{block}
\end{frame}

\begin{frame}{加权前后的平衡}
\begin{center}
\begin{tabular}{lrr}
\toprule
特征 & 原始 SMD & 加权 SMD\\
\midrule
risk score & 0.9710 & 0.2243\\
tenure months & -1.0255 & -0.3551\\
monthly spend & -1.0310 & -0.3600\\
support tickets & 1.0180 & 0.3257\\
usage sessions & -1.0033 & -0.3079\\
\bottomrule
\end{tabular}
\end{center}
\small
加权改善了可比性，但没有证明未观察混杂已经消失。
\end{frame}

\begin{frame}{差分中的差分}
\[
  \widehat{\tau}_{DID}
  =
  (\bar Y_{T,post}-\bar Y_{T,pre})
  -
  (\bar Y_{C,post}-\bar Y_{C,pre}).
\]
\begin{block}{平行趋势假设}
若没有干预，处理组和对照组的平均结果变化趋势应相同。
\end{block}
\small
DID 用对照组同期变化校正共同时间冲击，但不能自动处理不同趋势、溢出和同期政策变化。
\end{frame}

\begin{frame}{第8周 DID 数据}
\begin{center}
\begin{tabular}{lrrr}
\toprule
组别 & 干预前 & 干预后 & 变化\\
\midrule
处理门店 & 0.6100 & 0.7025 & 0.0925\\
对照门店 & 0.7125 & 0.7275 & 0.0150\\
\bottomrule
\end{tabular}
\end{center}
\[
  \widehat{\tau}_{DID}=0.0925-0.0150=0.0775.
\]
\small
在平行趋势假设下，留存策略使处理门店留存率额外提高约 7.75 个百分点。
\end{frame}

\begin{frame}[fragile]{代码入口}
\begin{lstlisting}[language=bash]
PYTHONPATH=src python3 src/bdm_decision/cases/week08_causality.py
\end{lstlisting}

\begin{lstlisting}[language=Python]
from bdm_decision.causal.propensity_score import (
    load_coupon_records, naive_difference, risk_band_adjusted_effect
)
from bdm_decision.causal.did import load_panel_records, difference_in_differences

records = load_coupon_records()
print(naive_difference(records))
print(risk_band_adjusted_effect(records))
print(difference_in_differences(load_panel_records()))
\end{lstlisting}
\end{frame}

\begin{frame}{课堂讨论}
\begin{enumerate}
  \item 为什么朴素差异为 0 仍不能说明优惠券无效？
  \item 哪些未观察因素可能同时影响发券和留存？
  \item 如果某类高风险客户全都被发券，重叠性会怎样失败？
  \item DID 的平行趋势可以用哪些业务证据支持？
  \item 预测模型、因果估计和优惠券预算应如何连接？
\end{enumerate}
\end{frame}

\begin{frame}{本讲小结}
\begin{itemize}
  \item 预测回答未来风险，因果回答行动效果；
  \item 潜在结果框架明确了反事实和处理效应；
  \item 随机实验是最清晰的识别基准；
  \item 观测性方法依赖一致性、SUTVA、可交换性和重叠性；
  \item 倾向得分改善可比性，但不能消除未观察混杂；
  \item DID 依赖平行趋势，应配合业务机制和前趋势检查。
\end{itemize}
\end{frame}

\end{document}
